\section{Categorical status}
\label{sec:categorical}

We work in $\PB^C_{\gauge_1}$ for a fixed $C$, and we assume
Convention~\ref{conv:degenerate}.

\subsection{Initial and terminal objects}

\begin{proposition}\label{prop:terminal}
$\RR$ is initial and the degenerate algebra $0$ is terminal.  Without
Convention~\ref{conv:degenerate}, and for $C\ge\tfrac14$, there is no terminal
object and hence no limits --- for reasons unrelated to the subject.
\end{proposition}

\begin{proof}
For initiality, the unit $\RR\to A$, $\lambda\mapsto\lambda1$, is the unique
unital homomorphism; it is isometric and both gauges vanish on its image.  For
terminality, $A\to0$ is unique, contractive and budget-contracting.

Without the degenerate object, suppose $T$ were terminal.  Since
$C\ge\tfrac14$, $M_2(\RR)$ is an object (Example~\ref{ex:M2}); let
$\varphi:M_2(\RR)\to T$ be the unique morphism.  As $M_2(\RR)$ is simple and
$\varphi$ is unital, $\varphi$ is injective.  For $v\in O(2)$ the automorphism
$\Ad_v$ of $M_2(\RR)$ is isometric, hence gauge-preserving by
Lemma~\ref{lem:isometric}, hence a PB-morphism; so $\varphi\circ\Ad_v$ is again
a morphism $M_2(\RR)\to T$, and uniqueness gives
$\varphi\circ\Ad_v=\varphi$.  Injectivity of $\varphi$ then forces
$\Ad_v=\id$ for every $v\in O(2)$, which is false.  Hence there is no terminal
object, and since the terminal object is the empty product, no limits.  (The
same argument runs with $\HH$ and its inner automorphisms in place of
$M_2(\RR)$; some noncommutative object is needed, which is why the statement is
restricted to $C\ge\tfrac14$.)
\end{proof}

\subsection{Products}

\begin{proposition}\label{prop:products}
$\PB^C_{\gauge_1}$ has all products.  For $A=\prod_iA_i$ with the sup norm and
$a=(a_i)$,
\[
  \gauge_1(a)=\sup_i\gauge_1(a_i),
\]
and the projections are PB-morphisms with the expected universal property.
\end{proposition}

\begin{proof}
For the gauge formula, $\norm{[a,z]}=\sup_i\norm{[a_i,z_i]}$ and $\norm{z}\le1$
forces $\norm{z_i}\le1$, giving $\le$; conversely a test element concentrated in
one factor gives $\ge$.

For \eqref{ax:pb1} let $M=\norm{a}=\sup_i\norm{a_i}$ and fix $\varepsilon>0$ and
an index $i$ with $\norm{a_i}\ge M-\varepsilon$.  Then
\[
  \norm{a}^2-\norm{a^2}
  \;\le\;M^2-\norm{a_i^2}
  \;=\;\bigl(M^2-\norm{a_i}^2\bigr)+\bigl(\norm{a_i}^2-\norm{a_i^2}\bigr)
  \;\le\;2M\varepsilon+C\gauge_1(a_i)
  \;\le\;2M\varepsilon+C\gauge_1(a),
\]
and $\varepsilon\downarrow0$ gives the claim.  For \eqref{ax:pb2} write
$S=\norm{a^2}=\sup_j\norm{a_j^2}$; then for each $i$,
\[
  \bigl\lVert S1-a_i^2\bigr\rVert
  \;\le\;\bigl\lVert\,\norm{a_i^2}1-a_i^2\,\bigr\rVert+\bigl(S-\norm{a_i^2}\bigr)
  \;\le\;\norm{a_i^2}+C\gauge_1(a_i)+S-\norm{a_i^2}
  \;\le\;S+C\gauge_1(a),
\]
and taking the supremum over $i$ on the left gives
$\bigl\lVert S1-a^2\bigr\rVert\le S+C\gauge_1(a)$.  Projections contract both norm and gauge; for
a cone $(f_i:T\to A_i)$ the map $f=(f_i)$ satisfies
$\gauge_1(f(x))=\sup_i\gauge_1(f_i(x))\le\gauge_1(x)$ and is contractive, and it
is the unique such factorisation.
\end{proof}

\begin{remark}[Uniformity of the constant is essential]
The proof uses one constant for all factors.  In $\PB_{\gauge_1}=\bigcup_C\PB^C$
products need not exist: a family with $C_{\gauge_1}(A_i)\to\infty$ has a product
whose critical constant is the supremum, hence infinite.  This is a reason to
regard the filtered pieces $\PB^C$, and not their union, as the categories of
interest.
\end{remark}

\subsection{Subalgebras, and why equalizers are hard}

\begin{proposition}[The class is not closed under subalgebras]
\label{prop:not-closed-subalgebras}
A closed unital subalgebra of a PB-algebra need not be a PB-algebra.  Indeed
$\RR1\oplus\RR E_{12}\cong\RR[\varepsilon]/(\varepsilon^2)$ is a closed unital
subalgebra of $M_2(\RR)$, and by Example~\ref{ex:excluded} it is not a
PB-algebra for any constant.
\end{proposition}

\begin{remark}[Ambient referencing]\label{rem:ambient-referencing}
Proposition~\ref{prop:not-closed-subalgebras} isolates the structural
difficulty.  The gauge is defined by a supremum over test elements $z$ drawn
from the ambient algebra, so passing to a subalgebra can destroy the budget
while preserving the defect that the budget was paying for.  A $C^*$-algebra
has no analogous failure: the datum defining the norm, $r(a^*a)^{1/2}$, is built
from $a$ and $a^*$, both of which lie in any $*$-subalgebra containing $a$.
Closure under $*$-subalgebras is free there, and that is exactly why $C^*$
has equalizers --- the solution set $\{f=g\}$ is a norm-closed $*$-subalgebra,
hence an object, and its inclusion is isometric, hence legal.  Everything
peculiar about limits in $\PB$ traces back to this one asymmetry.
\end{remark}

\begin{proposition}[The set-theoretic equalizer is not legal]
\label{prop:delta-not-saturated}
Let $A=M_2(\RR)$, $u=\operatorname{diag}(1,-1)$, and consider the parallel pair
$\id,\Ad_u:A\rightrightarrows A$.  The set-theoretic equalizer is the diagonal
$\Delta=\{\operatorname{diag}(\alpha,\beta)\}\cong\RR^2$, which \emph{is} a
PB-algebra, but its inclusion $\Delta\hookrightarrow M_2(\RR)$ is \emph{not} a
PB-morphism: $\gauge_1^\Delta\equiv0$ while
$\gauge_1^{M_2}(\operatorname{diag}(\alpha,\beta))=(\alpha-\beta)^2$, so
$\Delta$ is not saturated.  The largest saturated subalgebra of $\Delta$ is
$\RR1$.
\end{proposition}

\begin{proof}
$\Delta$ is commutative, so its own gauge vanishes.  In $M_2(\RR)$,
$[\operatorname{diag}(\alpha,\beta),z]$ has entries
$(\alpha-\beta)z_{12}$ and $-(\alpha-\beta)z_{21}$ off the diagonal and zero on
it, so its norm is $\abs{\alpha-\beta}\max(\abs{z_{12}},\abs{z_{21}})\le\abs{\alpha-\beta}$
with equality for $z=E_{12}$; this also agrees with
Proposition~\ref{prop:real-stampfli}, since
$\dist(\operatorname{diag}(\alpha,\beta),\RR1)=\tfrac12\abs{\alpha-\beta}$.  The
subalgebras of $\Delta\cong\RR^2$ are $\RR1$ and $\Delta$, and $\RR1$ is
saturated because both gauges vanish on it.
\end{proof}

\begin{remark}[A retraction, upheld]
The earlier notes first concluded from Proposition~\ref{prop:delta-not-saturated}
that $\PB$ is incomplete, then correctly retracted the argument: a categorical
equalizer need not be the set-theoretic one, and the natural candidate is
instead the largest saturated subalgebra $\RR1$, whose universality has to be
tested against \emph{legal} cones only.  That retraction stands.  What follows
settles the question it left behind.
\end{remark}

\subsection{Legal cones: the question, and its answer}

\begin{question}[\cite{notes}, Question~19]\label{q:legal-cone}
Does there exist an object $T$ and a PB-morphism $t:T\to M_2(\RR)$ with
$t(T)\subseteq\Delta$ and $t(T)\ne\RR1$?
\end{question}

The question has a clean reformulation, which also explains why it is not
obviously answerable either way.

\begin{lemma}[Cones are admissible character pairs]\label{lem:cones-characters}
PB-morphisms $t:T\to M_2(\RR)$ with $t(T)\subseteq\Delta$ correspond bijectively
to ordered pairs $(\chi_1,\chi_2)$ of real characters of $T$ satisfying
\begin{equation}\label{eq:admissible}
  \bigl(\chi_1(x)-\chi_2(x)\bigr)^2\;\le\;\gauge_1^T(x)\qquad\text{for all }x\in T,
\end{equation}
via $t(x)=\operatorname{diag}(\chi_1(x),\chi_2(x))$.  Every such $t$ equalizes
$\id$ and $\Ad_u$.  Call a pair satisfying \eqref{eq:admissible}
\emph{admissible}.
\end{lemma}

\begin{proof}
A unital homomorphism into $\Delta\cong\RR^2$ is a pair of unital
homomorphisms $T\to\RR$, i.e.\ a pair of real characters.  Contractivity is
automatic: $\norm{\operatorname{diag}(\chi_1(x),\chi_2(x))}=\max_i\abs{\chi_i(x)}\le r(x)\le\norm{x}$.
The budget condition is $\gauge_1^{M_2}(t(x))\le\gauge_1^T(x)$, and
$\gauge_1^{M_2}(\operatorname{diag}(\alpha,\beta))=(\alpha-\beta)^2$ by
Proposition~\ref{prop:delta-not-saturated}.  Finally $\Ad_u$ fixes $\Delta$
pointwise, so any such $t$ equalizes the pair.
\end{proof}

\begin{lemma}[Admissible pairs agree on the centre]\label{lem:agree-on-centre}
If $(\chi_1,\chi_2)$ is admissible then $\chi_1=\chi_2$ on $\Zc(T)$.
Equivalently, the two characters lie over the same point of the classical base
$X_T$.
\end{lemma}

\begin{proof}
For $x\in\Zc(T)$ we have $\gauge_1^T(x)=0$, so \eqref{eq:admissible} forces
$\chi_1(x)=\chi_2(x)$.
\end{proof}

So a legal cone with non-scalar image requires an object with two distinct
characters in a single fibre of its base --- a pointless-in-the-relevant-fibre
object that is noncommutative enough to fund the separation of its own
characters.  The commutative objects cannot do it, by
Lemma~\ref{lem:agree-on-centre}; $M_2(\RR)$ cannot do it, having no characters
at all.  But there is an object that can.

\begin{theorem}[Answer to Question~\ref{q:legal-cone}, for $C\ge\kap^*$]
\label{thm:t2-cone}
Let $C\ge\kap^*=0.3098\ldots$, so that $T=T_2(\RR)$ is an object of
$\PB^C_{\gauge_1}$ (Proposition~\ref{prop:T2-constant}), and let
\[
  t:T_2(\RR)\longrightarrow M_2(\RR),
  \qquad
  t\begin{pmatrix}\alpha&\gamma\\0&\beta\end{pmatrix}
  =\begin{pmatrix}\alpha&0\\0&\beta\end{pmatrix}
\]
be the map killing the radical.  Then $t$ is a PB-morphism, it equalizes $\id$
and $\Ad_u$, and its image is all of $\Delta$.  Consequently $\RR1$ is
\emph{not} the equalizer of $\id$ and $\Ad_u$ in $\PB^C_{\gauge_1}$.
\end{theorem}

\begin{proof}
$t$ is the quotient by $\rad(T_2(\RR))=\RR E_{12}$ followed by the identification
$T_2(\RR)/\RR E_{12}\cong\RR^2\cong\Delta$, so it is a unital homomorphism with
image $\Delta$.  Its two components are the characters
$\chi_1(x)=\alpha$ and $\chi_2(x)=\beta$; by Lemma~\ref{lem:cones-characters} it
remains to verify admissibility \eqref{eq:admissible}.

Let $x=\begin{psmallmatrix}\alpha&\gamma\\0&\beta\end{psmallmatrix}$ and
$z=\begin{psmallmatrix}p&c\\0&q\end{psmallmatrix}$.  A direct computation gives
\[
  [x,z]=\bigl(c(\alpha-\beta)+\gamma(q-p)\bigr)E_{12},
\]
so taking $z=E_{12}$, which lies in $T_2(\RR)$ and has norm $1$, yields
$\norm{[x,E_{12}]}=\abs{\alpha-\beta}$ and therefore
\[
  \gauge_1^{T_2}(x)\;\ge\;(\alpha-\beta)^2=\bigl(\chi_1(x)-\chi_2(x)\bigr)^2 .
\]
That is exactly \eqref{eq:admissible}.  Finally, if $\RR1$ with its inclusion
were the equalizer, the cone $t$ would factor through it, forcing
$t(T_2(\RR))\subseteq\RR1$, contradicting $t(T_2(\RR))=\Delta$.
\end{proof}

\begin{remark}[The mechanism]
$T_2(\RR)$ has two characters, the two diagonal entries, separated by a
non-central radical direction; that direction is precisely what funds the budget
realising their separation.  In the language of \S\ref{sec:visibility} it is an
object whose noncommutativity is Jacobson-invisible but commutator-visible ---
and it is the commutator that the morphism law measures.
\end{remark}

\begin{warning}[The answer depends on the constant]\label{warn:cone-constant}
Theorem~\ref{thm:t2-cone} requires $T_2(\RR)$ to be an \emph{object}, and by
Proposition~\ref{prop:T2-constant} it is one only for $C\ge\kap^*>\tfrac14$.  So
the status of Question~\ref{q:legal-cone} is not uniform in $C$:
\begin{itemize}
\item for $C\ge\kap^*$ the answer is \emph{yes} and $\RR1$ is not the equalizer;
\item for $C\in[\tfrac14,\kap^*)$ --- in particular at the critical constant
  $C=\tfrac14$, where $M_2(\RR)$ lives and $T_2(\RR)$ does not --- the question
  remains \emph{open}: the cone constructed above is unavailable, and we know of
  no other.
\end{itemize}
This is worth stating plainly because it is easy to miss.  A cone is an object
of the category, so an answer to a limit question can change as the constant
crosses a threshold, and $\kap^*$ is such a threshold.  Note also what
Lemma~\ref{lem:agree-on-centre} demands of any cone at $C=\tfrac14$: an object
with two distinct characters lying over one point of its base, with their
separation funded by its own commutators.  If
Conjecture~\ref{conj:semisimple} holds, such an object is semisimple, which
makes the search considerably harder --- the obvious source of two characters
in one fibre is a non-central radical, and that is exactly what
semisimplicity forbids.
\end{warning}

\begin{question}[Does the equalizer exist?]\label{q:equalizer-exists}
Theorem~\ref{thm:t2-cone} eliminates the natural candidate but does not decide
existence.  By Lemma~\ref{lem:cones-characters} the equalizer of $\id$ and
$\Ad_u$ exists if and only if the functor
\[
  F:\PB^C_{\gauge_1}\longrightarrow\cat{Set},
  \qquad
  F(T)=\bigl\{\text{admissible pairs of characters of }T\bigr\}
\]
is representable, in which case the representing object, equipped with its
universal admissible pair, is the equalizer.
\end{question}

This is a sharper question than the original, and it is the form in which we
recommend attacking completeness.  Two remarks on what a solution must handle.
$F(M_2(\RR))=\varnothing$, so a representing object $E$ admits no morphism from
$M_2(\RR)$; and $F$ takes a commutative $C(X,\RR)$ to the diagonal pairs only,
so $\Hom(C(X,\RR),E)$ must be the set of continuous maps $X_E\to X$ by
Corollary~\ref{cor:points} --- a strong constraint on the base of $E$.

\subsection{Colimits}

\begin{proposition}\label{prop:comm-colimits}
By Theorem~\ref{thm:coreflection} the inclusion
$\CommPB\hookrightarrow\PB^C_{\gauge_1}$ is a left adjoint, hence preserves all
colimits that exist in $\CommPB\simeq\CompHaus\op$.  In particular coproducts of
commutative objects exist and are given by products of bases
(Corollary~\ref{cor:comm-coproducts}).
\end{proposition}

\begin{remark}[Status of cocompleteness]
Warning~\ref{warn:not-free-product} disposes of the programme of the earlier
notes, which proposed to establish cocompleteness by checking that free products
of PB-algebras are PB-algebras: the free product is not the coproduct here.
Coequalizers reduce to the question of whether $\PB$ is closed under quotients
by closed two-sided ideals --- quotient maps are always legal by
Proposition~\ref{prop:category}, so the only issue is whether the quotient
satisfies the axioms --- which is Problem~\ref{prob:quotients}, and which is
also exactly what the bundle picture of \S\ref{sec:bundle} needs.  We regard
``cocomplete but not complete'' as the expected outcome, and note for the record
that the earlier notes' retraction of the reflector argument was correct: a
reflective subcategory of a complete category inherits limits, so a reflector
$\BanAlg\to\PB$ would force completeness.
\end{remark}
